\documentclass[11pt,letterpaper]{article}
\ifdefined\pdfminorversion\pdfminorversion=7\fi
\usepackage[T1]{fontenc}
\usepackage[utf8]{inputenc}
\usepackage{lmodern}
\usepackage[margin=0.88in,headheight=14pt]{geometry}
\usepackage{microtype}
\usepackage{amsmath,amssymb,amsthm,mathtools}
\usepackage{booktabs,tabularx,array}
\usepackage[dvipsnames]{xcolor}
\usepackage{listings}
\usepackage{enumitem,fancyhdr}
\usepackage[colorlinks=true,linkcolor=MidnightBlue,citecolor=MidnightBlue,
  urlcolor=MidnightBlue,breaklinks=true]{hyperref}
\hypersetup{pdftitle={Randomized Quantiles of Weighted Bernoulli Sums},
  pdfsubject={Exact randomization and conservative fallback in weightedBernoulli}}
\newtheorem{proposition}{Proposition}[section]
\newtheorem{remark}[proposition]{Remark}
\DeclareMathOperator{\supp}{supp}
\newcommand{\Pbb}{\mathbb P}
\newcommand{\E}{\mathbb E}
\definecolor{CodeBackground}{HTML}{F5F7FA}
\definecolor{CodeBorder}{HTML}{D9E0E8}
\lstdefinestyle{Rcode}{language=R,
  basicstyle=\ttfamily\fontsize{8.8}{10.6}\selectfont,
  keywordstyle=\color{MidnightBlue},commentstyle=\color{gray},
  stringstyle=\color{BrickRed},backgroundcolor=\color{CodeBackground},
  frame=single,rulecolor=\color{CodeBorder},framerule=0.3pt,
  framesep=5pt,aboveskip=8pt,belowskip=8pt,
  showstringspaces=false,upquote=true,keepspaces=true,
  columns=fullflexible,breaklines=true,breakatwhitespace=false,tabsize=2}
\lstset{style=Rcode}
\setlist{topsep=3pt,itemsep=2pt,parsep=0pt}
\setlength{\parskip}{4pt}
\setlength{\parindent}{0pt}
\setlength{\emergencystretch}{2em}
\pagestyle{fancy}
\fancyhf{}
\fancyhead[L]{\small Weighted Bernoulli quantiles}
\fancyhead[R]{\small Randomization addendum}
\fancyfoot[C]{\thepage}
\renewcommand{\headrulewidth}{0.3pt}
\title{\vspace{-1.1em}\textbf{Randomized Quantiles of\\Weighted Bernoulli Sums}\\[0.4em]
  \large An addendum to the computation note}
\author{}
\date{October 7, 2026}

\begin{document}
\maketitle
\vspace{-1.5em}
\begin{abstract}
For a discrete distribution, a deterministic quantile generally leaves
unused probability at its boundary. Independent randomization between two
adjacent support points removes this discrepancy when the distribution can
be computed exactly. This addendum gives the one-sided construction and the
original two-sided construction for fair weighted Bernoulli sums, proves
their probability identities, and explains the special treatment of an atom
at the center. It also specifies a conservative fallback and its different
output interpretation. Both constructions use the existing single public
function, \texttt{qweighted\_bernoulli()}.
\end{abstract}

\section{Model, interface, and probability conventions}

Let
\[
G=\sum_{i=1}^{B}w_i\xi_i,\qquad
\xi_i\sim\operatorname{Ber}(\pi_i)\text{ independently},\qquad
w_i\ge0,\quad 0\le\pi_i\le1,
\]
and write $F(x)=\Pbb(G\le x)$. The weights and success probabilities are
specified inputs. All probability identities involving a randomized cutoff
include a fresh uniform variable $V\sim U(0,1)$, independent of $G$ and of
any data or tie-breaking variables used elsewhere in the procedure.

The package retains one exported function. Three arguments are appended to
the existing interface, preserving earlier positional calls:
\begin{lstlisting}
qweighted_bernoulli(
  p, weights, prob = 0.5, max_length = 25,
  lower.tail = TRUE, log.p = FALSE, details = FALSE,
  control = list(),
  randomized = FALSE, two.sided = FALSE, draw = TRUE
)
\end{lstlisting}
The default \texttt{randomized = FALSE} retains deterministic quantile
computation. For \texttt{randomized = TRUE}, the function always returns a
list containing \texttt{c\_plus}, \texttt{c\_minus}, \texttt{tau},
\texttt{P\_plus}, \texttt{P\_minus}, \texttt{exact}, and \texttt{kappa}.
The field \texttt{tau} is always the probability of selecting the larger
cutoff \texttt{c\_plus}. Setting \texttt{draw = FALSE} returns the mixing
distribution with \texttt{kappa = NULL}; it does not generate a draw.

The argument \texttt{p} retains the usual probability convention. With
\texttt{lower.tail = FALSE}, it specifies the tail level $\alpha$ directly.
With \texttt{lower.tail = TRUE}, it specifies $1-\alpha$. Setting
\texttt{log.p = TRUE} supplies the logarithm of the corresponding probability.
The statements below use $0<\alpha<1$.

\section{Exact one-sided randomization}

Put $p_0=1-\alpha$ and define
\[
c_+=\min\{x\in\supp(G):F(x)\ge p_0\},\qquad
c_-=\max\{x\in\supp(G):x<c_+\}.
\]
If the latter set is empty, use the sentinel $c_-=-\infty$. This convention
also works when deterministic summands shift the minimum of the support
away from zero. Define
\begin{equation}\label{eq:one-sided-P}
P_+=\Pbb(G>c_+),\qquad P_-=\Pbb(G>c_-),\qquad
\tau=\frac{P_--\alpha}{P_--P_+}.
\end{equation}
Minimality of $c_+$ gives $P_+\le\alpha<P_-$ and hence $0<\tau\le1$.
The randomized cutoff is
\begin{equation}\label{eq:mixture}
\kappa=\begin{cases}
c_+,&V\le\tau,\\
c_-,&V>\tau.
\end{cases}
\end{equation}

\begin{proposition}\label{prop:one-sided}
The cutoff in \eqref{eq:mixture} satisfies
$\Pbb(G>\kappa)=\alpha$, or equivalently
$\Pbb(G\le\kappa)=1-\alpha$.
\end{proposition}
\begin{proof}
Independence permits conditioning on the selected cutoff, giving
\[
\Pbb(G>\kappa)=\tau P_++(1-\tau)P_-
 =P_--\tau(P_--P_+)=\alpha.
\]
The complementary identity follows because the event uses a strict upper
tail. No continuity assumption is needed.
\end{proof}

Equivalently,
\[
\tau=\frac{p_0-\Pbb(G<c_+)}{\Pbb(G=c_+)}.
\]
This identity explains the need for the atom at the quantile. Merely
returning the ordinary quantile does not provide the randomization. For a
constant $G=g_0$, the construction gives $c_+=g_0$, $c_-=-\infty$, and
$\tau=p_0$. The natural one-sided endpoint extensions are a cutoff of
$-\infty$ at $p_0=0$ and the maximum support point at $p_0=1$.

\begin{lstlisting}
one <- qweighted_bernoulli(
  p = 0.05, weights = c(1, 2, 4),
  prob = c(0.2, 0.5, 0.8), lower.tail = FALSE,
  randomized = TRUE, draw = FALSE
)
one[c("c_plus", "c_minus", "tau", "P_plus", "P_minus")]

# In the exact branch, this is the requested strict-tail level.
one$tau * one$P_plus + (1 - one$tau) * one$P_minus
\end{lstlisting}
Here \texttt{P\_plus} and \texttt{P\_minus} are the strict upper-tail
probabilities in \eqref{eq:one-sided-P}. In particular, the probability
associated with the larger cutoff is the smaller tail probability.

\section{The original fair two-sided cutoff}

Now assume $\pi_i=1/2$ for every positive weight, and let $N=\sum_iw_i$.
Complementing all Bernoulli variables gives $G\overset d=N-G$. Define
\[
\mathcal S=\{-1\}\cup\bigl(\supp(G)\cap[0,N/2]\bigr),
\]
and, for $c\in\mathcal S\setminus\{-1\}$,
\[
P(c)=\Pbb\!\left(\left|G-\frac N2\right|\ge\frac N2-c\right),
\qquad P(-1)=0.
\]
The value $-1$ is a sentinel below the nonnegative support. It is appropriate
even for noninteger weights: an increment of one between support points is
neither assumed nor used.

\begin{proposition}\label{prop:two-sided}
Let
\[
c_+=\min\{c\in\mathcal S:P(c)\ge\alpha\},\qquad
c_-=\max\{c\in\mathcal S:c<c_+\}.
\]
Then $c_+=q_G(\alpha/2)$, where $q_G$ is the ordinary leftmost distributional
quantile. Write $F_-=\Pbb(G<c_+)$. The probabilities at the two cutoffs are
\begin{equation}\label{eq:two-sided-P}
P_-=2F_-,\qquad
P_+=\begin{cases}
2F(c_+),&c_+<N/2,\\
1,&c_+=N/2.
\end{cases}
\end{equation}
With
\begin{equation}\label{eq:two-sided-tau}
\tau=\frac{\alpha-P_-}{P_+-P_-}
\end{equation}
and the independent mixture \eqref{eq:mixture},
\[
\Pbb\!\left(\left|G-\frac N2\right|\ge\frac N2-\kappa\right)=\alpha.
\]
\end{proposition}
\begin{proof}
For $c<N/2$, the events $\{G\le c\}$ and $\{G\ge N-c\}$ are disjoint.
Symmetry therefore gives $P(c)=2F(c)$. At $c=N/2$, their union is the whole
space, so $P(N/2)=1$. Symmetry also gives $F(N/2)\ge1/2$. Since
$\alpha/2<1/2$, the quantile $q_G(\alpha/2)$ is at most $N/2$ and is the first
support point at which $P(c)\ge\alpha$. Its predecessor is strictly below
$N/2$, and $F(c_-)=\Pbb(G<c_+)$, including the sentinel case. These facts prove
\eqref{eq:two-sided-P}. Minimality gives $P_-<\alpha\le P_+$, so the mixing
probability is in $(0,1]$. Conditioning on the independent mixture yields
$(1-\tau)P_-+\tau P_+=\alpha$, as claimed.
\end{proof}

\textbf{The center atom.} If $c_+=N/2$, the jump $P_+-P_-$ equals
$\Pbb(G=N/2)$, not twice that atom. Indeed, symmetry gives
$2F_-=1-\Pbb(G=N/2)$. Replacing $P_+$ by $2F(c_+)$ would double-count the
center and use the wrong randomization.

For $w=(1,1)$ and $\alpha=0.75$, the law of $G$ has masses $1/4,1/2,1/4$ at
$0,1,2$. Thus
\[
c_-=0,\quad c_+=1,\quad P_-=\frac12,\quad P_+=1,\quad\tau=\frac12.
\]
For all-zero weights, $G=N=0$ and the same definitions give $c_-=-1$,
$c_+=0$, and $\tau=\alpha$.

\begin{lstlisting}
center <- qweighted_bernoulli(
  p = 0.75, weights = c(1, 1), lower.tail = FALSE,
  randomized = TRUE, two.sided = TRUE, draw = FALSE
)
center[c("c_plus", "c_minus", "tau", "P_plus", "P_minus")]
# c_plus = 1, c_minus = 0, tau = 0.5
# P_plus = 1, P_minus = 0.5

cut <- qweighted_bernoulli(
  p = 0.05, weights = 1:20, lower.tail = FALSE,
  randomized = TRUE, two.sided = TRUE, draw = FALSE
)
cut$tau * cut$P_plus + (1 - cut$tau) * cut$P_minus
\end{lstlisting}
The two-sided mode uses the probabilities in \eqref{eq:two-sided-P}; the
larger cutoff now has the larger event probability. This explains the
different orientations of \eqref{eq:one-sided-P} and
\eqref{eq:two-sided-tau}, although \texttt{tau} selects the larger cutoff in
both modes.

\section{Obtaining the bracket without full enumeration}

The exact engines and the rule \texttt{max\_length = 25} are described in
the main note~\cite{mainnote}. Randomization uses the same exact law and adds
the predecessor and the boundary probabilities. It does not require a new
public function or full enumeration of all $2^B$ Bernoulli outcomes.

In a meet-in-the-middle calculation, write $G=A+D$, where $A$ and $D$ are
independent half-sums with sorted supports $a_i$ and $d_j$ and masses
$r_i$ and $s_j$. At a candidate $x$, determine
\[
j_i^{\le}(x)=\max\{j:a_i+d_j\le x\},\qquad
j_i^{<}(x)=\max\{j:a_i+d_j<x\},
\]
using zero when the respective set is empty. With prefix masses
$S_j=\sum_{k\le j}s_k$ and $S_0=0$, the exact queries are
\[
F(x)=\sum_i r_iS_{j_i^{\le}(x)},\qquad
\Pbb(G<x)=\sum_i r_iS_{j_i^{<}(x)}.
\]
The actual predecessor of $x$ is
\[
\max_{i:j_i^{<}(x)>0}\bigl(a_i+d_{j_i^{<}(x)}\bigr),
\]
or the appropriate sentinel if the index set is empty. These same strict
and non-strict queries already verify the jump containing an exact
quantile. Reusing them avoids enumerating the Cartesian product of the
two half-supports.

The repeated-class engine applies the same calculation after replacing
equal weight--probability pairs by binomial counts. Binomial and Wilcoxon
supports have consecutive lattice coordinates. A general lattice
distribution can have gaps, so its predecessor must be the previous
reachable coordinate. Subtracting a unit or an arbitrary numerical tolerance
from an exact quantile is not a general substitute for this search.

\textbf{Numerical meaning of exact.} The probability identities above are
mathematical identities. As in the main note, \texttt{exact = TRUE} describes
an exact finite-law algorithm implemented in double precision; it does not
assert interval-certified arithmetic. Direct small-tail evaluation,
logarithmic probabilities, and stable differences avoid unnecessary
cancellation. On a recognized lattice, integer coordinates also determine
whether the quantile is the center. Arbitrary tolerance-based merging of
nearby atoms is not justified by the theory.

\section{A conservative fallback with a stated guarantee}

If the exact law is unaffordable, the existing algorithm computes the
applicable conservative candidates under its work limits and chooses their
minimum~\cite{mainnote}. The selected deterministic bound is chosen before
any random draw. A bound alone does not identify the quantile atom, and its
upper-tail certificate cannot be inserted into an exact interpolation
formula as though it were the true tail probability.

\subsection*{One-sided mode}

Suppose the selected threshold $q$ has a certificate
\[
\Pbb(G>q)\le b\le\alpha.
\]
The fallback reports the degenerate mixture
$c_+=c_-=q$ and $\tau=1$. It consequently satisfies
$\Pbb(G>\kappa)\le\alpha$. The fields \texttt{P\_plus} and
\texttt{P\_minus} report bounds on the event probabilities, rather than
the event probabilities themselves, and \texttt{exact = FALSE} identifies
this case. The two reported cutoffs do not assert that the true quantile has
been bracketed by adjacent support points.

\subsection*{The original two-sided mode}

Suppose the selected upper quantile bound at level $\alpha/2$ has a certificate
\begin{equation}\label{eq:qcertificate}
\Pbb(G>q)\le b\le\alpha/2.
\end{equation}
Choose an actual support value $c$ satisfying
\begin{equation}\label{eq:strict-support}
0\le c\le N/2,\qquad c<N-q,
\end{equation}
or choose the sentinel $c=-1$ if no such value is obtained.

\begin{proposition}\label{prop:conservative}
Every cutoff chosen in this way satisfies $P(c)\le\alpha$.
\end{proposition}
\begin{proof}
The sentinel has $P(-1)=0$. Otherwise, symmetry and the strict inequality in
\eqref{eq:strict-support} imply
\[
F(c)\le\Pbb(G<N-q)=\Pbb(N-G>q)=\Pbb(G>q)\le b.
\]
In fact $c<N/2$: if $c=N/2$, then $F(c)\ge1/2$ by symmetry, contradicting
$b\le\alpha/2<1/2$. Thus $P(c)=2F(c)\le2b\le\alpha$.
\end{proof}

A deterministic greedy subset search gives an inexpensive support value.
Start at zero when it satisfies \eqref{eq:strict-support}; traverse the
positive weights and add a weight only when the resulting sum still
satisfies that condition. Run the search in ascending and descending order
and retain the larger feasible result. Each result is a sum of a subset of
the original weights, so it belongs to $\supp(G)$. Optimal subset packing
is unnecessary for the guarantee: the strict budget and actual support
membership suffice. The sentinel remains available when the budget is
nonpositive. Numerical comparisons retain the precision qualification
stated above.

The fallback reports $c_+=c_-=c$, $\tau=1$, and
\texttt{exact = FALSE}. Its \texttt{P\_plus} and \texttt{P\_minus} are
certified bounds, at most $2b$ (zero for the sentinel). These equal endpoints
are a degenerate conservative cutoff, not the original exact bracket.
The algorithm does not claim equality with $\alpha$ in this branch.

The strict condition $c<N-q$ is essential. A strict upper-tail bound controls
$\Pbb(G<N-q)$ by symmetry; setting $c=N-q$ can add a boundary atom that the
certificate does not cover.

\section{Drawing and interpreting the result}

To obtain the mixing distribution without consuming a random draw, use
\texttt{draw = FALSE}. To obtain the draw as well, use the default
\texttt{draw = TRUE}:
\begin{lstlisting}
set.seed(20261007)
out <- qweighted_bernoulli(
  p = 0.05, weights = 1:20, lower.tail = FALSE,
  randomized = TRUE, two.sided = TRUE
)
out$kappa
out$exact

# The returned mixing distribution can also be sampled explicitly.
law <- qweighted_bernoulli(
  p = 0.05, weights = 1:20, lower.tail = FALSE,
  randomized = TRUE, two.sided = TRUE, draw = FALSE
)
v <- runif(length(law$tau))
kappa <- ifelse(log(v) < law$log_tau, law$c_plus, law$c_minus)
\end{lstlisting}
The explicit draw must be independent of the variables to which the cutoff
will be applied. In particular, a tie-breaking uniform from an earlier
construction must not be reused. The mathematical assertions concern the
joint probability over $G$ and this independent randomization; conditional
on a realized cutoff, its individual event probability need not equal
$\alpha$.

R's pseudorandom generators have finite resolution; most supplied uniform
generators produce at most $2^{32}$ distinct values~\cite{rrng}. Thus the
ideal-uniform probability identities are subject to both arithmetic and
simulation discretization. In particular, a single uniform comparison
cannot realize every extremely small mixing probability exactly. For such
levels, \texttt{draw = FALSE} preserves the analytical mixture and its
logarithmic weights for a separate sampling scheme. Standard R seeding
controls reproducibility; the function does not reset the seed.

Check \texttt{exact} before interpreting the probability fields. In an exact
row, the weighted average of \texttt{P\_plus} and \texttt{P\_minus} equals
the requested level, subject to ordinary numerical precision. In a
conservative row, those fields are bounds and the weighted average certifies
an inequality. All candidate selection is deterministic in the supplied
weights, probabilities, levels, and work limits. Taking the minimum of
separately randomized upper thresholds is not covered by the guarantee:
even two individually valid randomized thresholds can produce an invalid
minimum.

\begin{thebibliography}{9}
\bibitem{mainnote}
\emph{Computing Quantiles of Weighted Bernoulli Sums: Exact Algorithms,
Conservative Bounds, and Reproducible R Code} (2026).
Technical note accompanying the \texttt{weightedBernoulli} package,
October 7, 2026. The exact computation engines, concentration inequalities,
work limits, and numerical comparisons are documented there. The note is
included as \texttt{weighted\_bernoulli\_quantiles.pdf} in the package
documentation. This addendum supplies the integrated randomized interface
and the probability identities needed to interpret it.
\bibitem{rrng}
R Core Team. \emph{Random Number Generation} and \emph{The Uniform Distribution}.
R documentation, accessed October 7, 2026.
\url{https://stat.ethz.ch/R-manual/R-devel/library/base/html/Random.html};
\url{https://stat.ethz.ch/R-manual/R-devel/library/stats/html/Uniform.html}.
\end{thebibliography}
\end{document}
